\chapter{Discussion and Conclusions}
\label{ch:conclusion}

This chapter consolidates the work of the previous chapters: what was built, what Chapter~\ref{ch:results} establishes and what it does not, and which directions are most worth pursuing next.

\section{Summary of contributions}
\label{sec:conclusion-contributions}

We produce two main contributions. The first is the process-quality score $J$, a formula that measures refactoring quality by combining the endpoint cleanliness gain with five penalty terms from the literature -- broken builds, skipped tests, manual edits when IDE refactoring would work, long gaps between commits, and an intermediate-cleanliness lag penalty that costs trajectories which linger below their final cleanliness target -- and a step-savings bonus that rewards alternatives reaching the same end state in fewer steps. The second is the divergence-point detector, which finds points in a refactoring session where a better alternative exists, classifies them into four kinds (\emph{IDE-Replay}, \emph{Rework}, \emph{Hygiene}, \emph{Ordering}), and produces a concrete alternative trajectory for each. The alternatives are generated by four per-kind synthesisers internal to the detector: a reorder enumerator that finds valid step reorderings and checks them for AST equivalence, an IDE replayer that converts manual edits into IDE refactor calls, a rework stripper that removes self-cancelling step pairs, and a hygiene synthesiser that models what would happen if tests were run or commits made at different points.

Supporting these are: a 45-session hand-recorded dataset on a purpose-built Java fixture with multi-label ground-truth annotations; a 12-session user study and a 6-session agent comparison on a separate fixture; the IntelliJ plugin that captures the event traces; the JCEF dashboard that surfaces the divergence points; and a reproducible notebook from which every number in Chapter~\ref{ch:results} can be regenerated.

\section{Headline findings}
\label{sec:conclusion-findings}

Five findings carry our main claims.

First, the lag term $W_{\textsc{lag}}\,\ell(\tau)$ is the operative lever for the \emph{Ordering} divergence-kind. Across the $45$ injection sessions the reorder synthesiser admitted $194$ permutations over $24$ windows; $10$ of those windows produced a positive-magnitude alternative ($41.7\%$), and a further $2$ produced a negative-magnitude alternative (the user front-loaded cleanliness better than the synthesised alts). On the five sessions whose playbook author intended an ordering they believed was worse than its alternatives, $1$ of $5$ produced a strictly better permutation. Without the lag term in the formula, the corresponding numbers collapse to $4 / 24$ positive-magnitude windows, $0$ negative-magnitude verdicts, and $0$ ordering injections caught -- the lift is mechanically attributable to $W_{\textsc{lag}}$, since the rest of the formula was held fixed across the comparison (\S\ref{sec:results-headline}, \S\ref{sec:results-divergence}).

Second, the kind classifier is precise on the controlled-injection sessions. Per-kind precision is $1.00$ across all four kinds against the inter-rater-agreed \texttt{expected\_kinds} labels, with Cohen's $\kappa = 1.00$ for \emph{Ordering} and \emph{IDE-Replay} on all three rater pairs and $\kappa = 0.72$--$1.00$ for \emph{Rework} and \emph{Hygiene} (\S\ref{sec:results-kappa}, \S\ref{sec:results-divergence}). Every injection caught with a positive-magnitude alternative trajectory ranked at top-$1$ within its session.

Third, the score formula sits in a locally stable region around its production weights. Top-$1$ stability under single-knob perturbation is $96.1\%$ on the user-study sessions. The multi-knob Monte Carlo design (every weight perturbed simultaneously from $\exp(\sigma \cdot \mathcal{N}(0, 1))$ with $\sigma = \ln 2$) lowers this to $81.3\%$ (\S\ref{sec:results-sensitivity}). The formula is robust to plausibly-sized single-weight changes and to a majority of plausibly-sized simultaneous changes, and degrades smoothly under larger ones.

Fourth, both human participants showed near-monotonic improvement in their process discipline across the user-study arc. Decomposing the process score with the cleanliness-gain weight zeroed isolates terms that do not depend on codebase familiarity. P1's score rises from $25$ at S1 to $50$ at S6 despite the tasks getting progressively harder; P2 trends in the same direction less consistently (\S\ref{sec:results-userstudy}). The agent's gain-stripped score doesn't improve over time; it varies by task difficulty, not session number.

Fifth, the agent (Claude Code, model \texttt{claude-opus-4-7}, tested on 2026-05-21) course-corrected between sessions on the procedural-discipline behaviours within its tool surface, and correctly identified the one behaviour that lay outside it. Three threads from the transcript (\S\ref{sec:results-userstudy-agent}) show this: commit cadence was revised twice (under-committing in S1, over-correcting in S3, resolved by S6 atomic-edit synthesis); smell-introduction self-diagnosis after S3 prevented recurrence in S4--S5; IDE-replay flags persisted across the arc not because the agent failed to read the feedback but because it edits files via text tool calls and has no IDE refactor surface available to it, a limitation it diagnoses correctly in its own S6 reflection.

\section{Discussion}
\label{sec:conclusion-scope}

The chapter substantiates four claims: a divergence-point classifier whose precision and rank-1 prominence are conditional on the recorded sessions but reproducible; a score formula whose ranking is locally robust to weight choices, with each term carrying measurable influence; a descriptive observation that two participants showed gain-stripped process-discipline improvement across the six-session arc after receiving feedback; and an observation that a coding agent adjusted its behaviour between sessions in response to the same written feedback, within the limits of its tool surface.

A few stronger claims are left for follow-up work. The user study and agent comparison are observational rather than controlled, so the process-discipline arc is consistent with the feedback driving it but is not reported here as a causal finding. The plugin and synthesisers are written against IntelliJ and JDT, so generalisation to other IDEs or languages is a port rather than a configuration change. The weights are grounded in literature-supported severity orderings rather than fit against an external outcome dataset, since no dataset linking session-level process signals to downstream maintenance outcomes exists at the granularity the score operates on. Chapter~\ref{ch:threats} discusses each of these in detail.

\section{Future work}
\label{sec:conclusion-future-work}

The suggestions below are ordered by what would most strengthen the main findings, with the scale of the human study the largest single gap:

\begin{enumerate}
    \item \textbf{Larger and controlled human study.} The strongest descriptive claim in the chapter (the gain-stripped process-discipline arc in \S\ref{sec:results-userstudy}) rests on two participants without a control condition. A study with more participants and a feedback-on / feedback-off arm would establish causation rather than just describe the effect, and would enable better estimation of between-participant variance.

    \item \textbf{Multi-agent comparison.} The agent-comparison chapter rests on one agent run with one model version on one date. Extending it to multiple models across multiple dates would separate model-specific behaviour from agent-class behaviour and would let the qualitative threads in \S\ref{sec:results-userstudy-agent} become quantitative.

    \item \textbf{Rater-rank validation.} The Cohen's $\kappa$ study validates kind-presence agreement only. A Kendall $\tau$-b study on per-session divergence-point rankings would validate that the score's ordering of divergence points within a session matches a human expert's ranking, which the $\kappa$ figures alone do not establish.

    \item \textbf{IDE-refactor-as-agent-tool deployment.} The IDE-replay structural limitation in \S\ref{sec:results-userstudy-agent} is the cleanest example in the chapter of feedback being correctly interpreted but inactionable: the agent acknowledges the IDE refactoring would score better but has no way to invoke it. Exposing IntelliJ's refactor actions as agent tool calls would close that gap and would let the \emph{IDE-Replay} count actually respond to feedback across sessions.
\end{enumerate}

\section{Closing}
\label{sec:conclusion-closing}

The introduction opened by noting that the literature evaluates refactoring as a mapping from a start state to an end state, with the trajectory between them largely outside the picture. Our work shifts the question being asked: from whether the code is better after the refactoring to whether the process the developer followed is one a careful developer would have followed. The point underneath the numbers in Chapter~\ref{ch:results} is that a refactoring trajectory has concrete alternative paths at identifiable points, and that the gap between the observed path and those alternatives is not something a snapshot-only evaluation can see.
